% META: {"id": "1NSI-WEB-IHM-RE-C7-CORRIGE", "chapitre": "1NSI-WEB-IHM", "type_objet": "corrige", "exercice_ref": "1NSI-WEB-IHM-RE-C7", "capacites": ["P-IHM-04A"], "capacites_codes": ["C7"], "mode_creation": "ex_nihilo", "sources_inspiration": [], "status": "needs_review"}
\begin{corrige}{1NSI-WEB-IHM-RE-C7}
\begin{enumerate}
  \item « Adresse électronique » : c'est le contenu de la balise \lstinline{<label>}. Ni
        \lstinline{id} ni \lstinline{name} n'apparaissent à l'écran.
  \item Sous le nom \lstinline{courriel}, valeur de l'attribut \lstinline{name} — et non
        \lstinline{c1}, qui est l'\lstinline{id} servant au code et au \lstinline{<label>}.
  \item \textbf{Non}, il n'a pas de \lstinline{name} : sa valeur n'est pas transmise. Il faut
        lui en donner un, par exemple
        \lstinline{<input type="text" id="c2" name="pseudo">}.
  \item L'attribut \lstinline{required}.
\end{enumerate}
\end{corrige}
